\chapter{H. pylori Urea Breath Test}

\section*{Overview}
The H. pylori urea breath test detects active infection with Helicobacter pylori by measuring labeled carbon dioxide in exhaled breath after a urea drink.

\section*{Why It Is Done}
It helps diagnose H. pylori infection and confirm eradication after treatment in appropriate patients with dyspepsia, peptic ulcer disease, or related risk.

\section*{How It Is Performed}
After required fasting and medicine restrictions, the patient provides a baseline breath sample, drinks labeled urea, and provides one or more additional breath samples for laboratory analysis.

\section*{Preparation and Results}
Recent antibiotics, bismuth, and acid-suppressing medicines can cause false-negative results and may need to be withheld under clinical direction. The result is reported as positive or negative for active infection.

\section*{References}
\begin{enumerate}
  \item MedlinePlus. Helicobacter Pylori Tests. U.S. National Library of Medicine; 2025.
  \item Current Medical Diagnosis and Treatment. McGraw-Hill; 2025.
\end{enumerate}
\clearpage
